\subsection{自由群}

设 \(G_x = Z\) ，对所有 \(x \in X\) 。族 \((G_x)_{x \in X}\) 的自由积称为在 \(X\) 上构造的自由群，记作 \(F(X)\) 。设 \(\phi_x\) 是 \(G_x = Z\) 到 \(F(X)\) 中的典范同态。根据第 3 号命题 5 的推论，映射 \(x \mapsto \phi_x(1)\) （从 \(X\) 到 \(F(X)\) ）是单射；我们将把 \(X\) 与它在该映射下在 \(F(X)\) 中的像等同起来。于是 \(X\) 生成 \(F(X)\) 且 \(e \notin X\) 。

应用第3号命题5，我们得到以下结果：

命题7. 设 \(g\) 是自由群 \(F(X)\) 的一个元素。存在一个整数 \(n \geqslant 0\) 以及一个序列 \((x_{\alpha}, m(\alpha))_{1 \leqslant \alpha \leqslant n}\) ，它由关系 \(x_{\alpha} \in X\) ,

\[
g = \prod_ {\alpha = 1} ^ {n} x _ {\alpha} ^ {m (\alpha)}
\]

唯一确定。自由群 \(F(X)\) 具有如下泛性质：

命题8. 设G是一个群，f是将X映入G的一个映射。存在且仅存在一个同态 \(\bar{f}\) 将F(X)映入G，且扩展 \(f\) .

唯一性 \(\bar{f}\) 由下述事实得出：群 \(\mathrm{F}(\mathbf{X})\) 由 \(\mathbf{X}\) 生成。对所有 \(x\) （属于 \(\mathbf{X}\) ），设 \(f_x\) 是同态 \(n \mapsto f(x)^n\) ，它把 \(\mathbf{Z}\) 映入 \(\mathbf{G}\) 。根据第 3 号命题 4，存在一个同态 \(\bar{f}\) ，它把 \(\mathrm{F}(\mathbf{X})\) 映入 \(\mathbf{G}\) ，使得 \(\bar{f}(x^n) = f_x(n)\) 对 \(x \in \mathbf{X}\) 和 \(n \in \mathbf{Z}\) 成立；特别地， \(\bar{f}(x) = f_x(1) = f(x)\) 对所有 \(x \in \mathbf{X}\) 成立，因此 \(\bar{f}\) 扩张 \(f\) 。

设 \(u: X \to Y\) 是一个映射。根据命题 8，存在且仅存在一个从 \(F(X)\) 到 \(F(Y)\) 的同态，它与 \(u\) 在 \(X\) 上一致；把它记作 \(F(u)\) 。如同广群的情形（第 1 号）一样，公式

\[
\mathbf {F} (v \circ u) = \mathbf {F} (v) \circ \mathbf {F} (u)
\]

对每个映射 \(v: \mathbf{Y} \to \mathbf{Z}\) 成立，并且证明了 \(\mathrm{F}(u)\) 在 \(u\) 为单射（相应地为满射、双射）时也为单射（相应地为满射、双射）。对每个子集 \(S\) （属于 \(X\) ）， \(F(S)\) 将与 \(F(X)\) 中由 \(S\) 生成的子群等同起来。

设 I 为一个集合。在某些情形下，不把 I 中的 i 与其典范像 \(\phi_i(1)\) （在自由群 \(\mathbf{F}(\mathbf{I})\) 中）等同起来是有意义的；后者将记作 \(\mathbf{T}_i\) （或 \(\mathbf{T}_i', \mathbf{X}_i, \ldots\) ，视情况而定），并称为指标 \(i\) 的不定元。自由群 \(\mathbf{F}(\mathbf{I})\) 随后记作 \(\mathbf{F}((\mathbf{T}_i)_{i \in 1})\) 或 \(\mathbf{F}(\mathbf{T}_1, \ldots, \mathbf{T}_n)\) （若 \(\mathbf{I} = \{1, 2, \ldots, n\}\) ）。

设 G 为一个群，且 \(\mathbf{t} = (t_{i})_{i \in \mathbf{I}}\) 是 G 的一个元素族。根据命题 8，存在一个同态 \(f_{t}\) ，它把 \(\mathrm{F}((\mathrm{T}_{i})_{i \in \mathbf{I}})\) 映入 G，其特征是 \(f_{\mathbf{t}}(\mathrm{T}_{i}) = t_{i}\) 对所有 \(i \in I\) 成立。 \(\mathrm{F}((\mathrm{T}_{i})_{i \in \mathbf{I}})\) 中的元素 w 在 \(f_{t}\) 下的像将记作 \(w(\mathbf{t})\) 或 \(w(t_{1}, \ldots, t_{n})\) （若 \(I = \{1, 2, \ldots, n\}\) ）；称 \(w(\mathbf{t})\) 是在 w 中以替换 \(T_{i} \mapsto t_{i}\) 所得的结果。特别地，若取 \(\mathrm{G} = \mathrm{F}((\mathrm{T}_{i})_{i \in \mathbf{I}})\) 且 \((t_{i}) = (\mathrm{T}_{i}) = \mathbf{T}\) ，则 \(f_{T}\) 是 G 的恒等同态，从而 \(w(\mathbf{T}) = w\) ；对 \(I = \{1, 2, \ldots, n\}\) ，则 \(w(\mathrm{T}_{1}, \ldots, \mathrm{T}_{n}) = w\) 。

设 G 与 G' 为两个群，u 为 G 到 G' 中的一个同态，且 \(\mathbf{t} = (t_{1}, \ldots, t_{n})\) 为 G 中元素的一个有限序列。设 \(\mathbf{t}' = (u(t_{1}), \ldots, u(t_{n}))\) ；同态 \(u \circ f_{t}\) （从 \(\mathrm{F}(\mathrm{T}_{1}, \ldots, \mathrm{T}_{n})\) 到 G' 中）把 \(T_{i}\) 映到 \(u(t_{i})\) （对 \(1 \leqslant i \leqslant n\) ），因此等于 \(f_{t'}\) ；对 \(\mathrm{F}(\mathrm{T}_{1}, \ldots, \mathrm{T}_{n})\) 中的 w，则

\[
u (w (t _ {1}, \dots , t _ {n})) = w (u (t _ {1}), \dots , u (t _ {n})).\tag{8}
\]

设给定了 \(w\) （在 \(F(T_1, \ldots, T_n)\) 中）以及元素 \(v_1, \ldots, v_n\) （在自由群 \(F(T_1', \ldots, T_m')\) 中）。替换 \(T_i \mapsto v_i\) 定义了元素 \(w' = w(v_1, \ldots, v_n)\) （属于 \(F(T_1', \ldots, T_m')\) ）。设 \(G\) 为一个群， \(t_1, \ldots, t_m\) 为 \(G\) 的元素，而 \(u\) 为 \(F(T_1', \ldots, T_m')\) 到 \(G\) 中的同态，其特征是 \(u(T_j') = t_j\) 对 \(1 \leqslant j \leqslant m\) 成立。则 \(u(v_i) = v_i(t_1, \ldots, t_m)\) 且 \(u(w') = w(t_1, \ldots, t_m)\) ；公式 (8) 由此蕴含

\[
w ^ {\prime} (t _ {1}, \dots , t _ {m}) = w (v _ {1} (t _ {1}, \dots , t _ {m}), \dots , v _ {n} (t _ {1}, \dots , t _ {m})).\tag{9}
\]

这证明了“函数记号” \(w(t_{1},\ldots,t_{n})\) 的合理性。读者需自行将公式 (8) 和 (9) 推广到任意指标集的情形。

